\section{Resumo}
\label{sec:resumo}
%
Este projeto consiste no desenvolvimento de um robô autónomo capaz de detetar obstáculos e afastar-se deles, à medida que se aproxima do alvo pretendido. Para a criação deste robô, foi necessário seguir uma determinada ordem de trabalhos para obter o produto final. 

Primeiramente, foi necessário o desenvolvimento do modelo do robô, recorrendo-se a uma ferramenta CAD (Autodesk Fusion 360). Nesta plataforma, foram concebidas as diferentes peças que constituem o modelo do robô.  

De seguida, todas as peças em CAD foram importadas da plataforma para o CoppeliaSim, de forma a criar o modelo físico do robô.  

Por fim, para a navegação autónoma, foram implementados os algoritmos baseados em sistemas dinâmicos não lineares em Matlab e também foi realizada a comunicação entre o Matlab e o CoppeliaSim para o controlo do robô.